\documentclass[11pt]{article}
\usepackage[margin=1in]{geometry}
\usepackage[T1]{fontenc}
\usepackage[utf8]{inputenc}
\usepackage{lmodern}
\usepackage{setspace}
\usepackage{amsmath,amssymb,mathtools}
\usepackage{booktabs}
\usepackage{graphicx}
\usepackage{float}
\usepackage{hyperref}
\usepackage{caption}
\usepackage{pgfplotstable}
\pgfplotsset{compat=1.18}

\title{Methods: Polymarket Sentiment-to-Prediction Pipeline\\\large Teichmann-Style Randomized Signatures}
\author{Methodology Document}
\date{\today}

\begin{document}
\maketitle
\onehalfspacing

\subsection{Objective}
This methodology estimates whether exogenous sentiment can improve out-of-sample prediction of Polymarket forward returns under realistic transaction costs. The pipeline consists of topic-conditioned market selection, synchronized sentiment--price path construction, randomized-signature feature extraction through controlled differential equations, ridge-based statistical learning, and fee-stressed backtesting.
\paragraph{Research question in plain language.}
The practical question is straightforward: when public discourse becomes more positive or negative about a geopolitical theme, does that information appear early enough in sentiment streams to help forecast short-horizon changes in prediction-market prices. The methodological design therefore treats sentiment as an external information channel and evaluates whether that channel contributes predictive content beyond lagged price alone.
\paragraph{Design principle.}
The document follows a conservative empirical philosophy: model complexity is allowed in the feature map, but evaluation is constrained by out-of-sample testing, walk-forward retraining, non-overlapping horizons, and explicit transaction costs. This combination is intended to reduce the gap between statistical significance and tradable relevance.

\subsection{Notation and data model}
Let \(t_0 < t_1 < \cdots < t_T\) be a regular grid after resampling at interval \(\Delta t\).
Define \(S_t \in \mathbb{R}\) as standardized sentiment, \(P_t \in (0,1)\) as market price/probability clipped to \([\varepsilon,1-\varepsilon]\), \(\ell\) as lag in grid steps, \(h\) as forecast horizon in grid steps, and \(w\) as rolling window length.

For each decision time \(t\), define input path window
\[
X^{(t)} = \left\{(S_u, P_{u-\ell})\right\}_{u=t-w}^{t} \in (\mathbb{R}^2)^{w+1},
\]
and forward-return target
\[
r_t^{(h)} = \log P_{t+h} - \log P_t.
\]

\paragraph{Assumption A1 (Synchronous observability).}
After alignment and fill policy, \(S_t\) and \(P_t\) are both available on the same grid.

\paragraph{Assumption A2 (Piecewise-linear interpolation).}
Within each grid interval, the path is treated as piecewise linear, enabling discrete CDE integration.
\paragraph{Why this representation is used.}
The path formulation \((S_t, P_{t-\ell})\) is chosen because it preserves temporal ordering and interaction structure. Rather than forcing handcrafted nonlinear transforms, the path is passed to signature-based machinery that is specifically designed to capture sequential effects, including lead--lag and interaction geometry, in a mathematically controlled way.

\subsection{Market selection protocol}
Given keyword set \(\mathcal{K}\) and date window \([T_{\min}, T_{\max}]\), candidate markets are queried from Gamma and retained only when the market question matches the keyword regex with word boundaries, the contract is binary with valid CLOB token IDs, and non-zero price-history overlap is observed in \([T_{\min}, T_{\max}]\). Candidates are ranked by in-window price-point count. This is implemented in \texttt{scan\_overlap\_markets.py}.
\paragraph{Identification logic.}
Keyword filtering and overlap constraints are not only engineering conveniences; they define the empirical estimand. The model is not asked to learn a universal mapping across unrelated topics. Instead, it is evaluated in topic-consistent windows where sentiment plausibly carries domain information about the selected contracts.

\subsection{Randomized-signature feature construction}
\paragraph{Controlled differential equation.}
For path \(X\in C([0,T],\mathbb{R}^d)\), \(d=2\), hidden state \(Y\in\mathbb{R}^N\):
\[
Y_t = y_0 + \sum_{i=1}^{d} \int_0^t V_i(Y_s)\, dX_s^i.
\]

\paragraph{Vector field families.}
Two families are used:
\begin{align*}
\text{Depth-1:}\quad &V_i(y) = \tanh(A_i y + b_i),\\
\text{Depth-2 exp:}\quad &V_i(y) = \exp\!\left(A_i\,\exp(D_i y)\right),
\end{align*}
with random \(A_i, b_i, D_i\).
The depth-2 exponential form follows the Teichmann-style construction used for stronger signature reconstruction behavior.
\paragraph{Intuition.}
The CDE can be read as a continuous-time recurrent system whose hidden state reacts to each component of the path. Randomized vector fields generate a diverse basis of nonlinear path responses. The resulting terminal states are not interpreted individually; they are used collectively as a high-capacity, yet linear-regression-friendly representation.

\paragraph{Feature map.}
For each window \(X^{(t)}\), solve the CDE for \(M\) random initial states \(\{y_0^{(j)}\}_{j=1}^M\):
\[
\Phi(X^{(t)}) =
\left[
Y_T^{(1)}(X^{(t)}),\dots,Y_T^{(M)}(X^{(t)})
\right]\in \mathbb{R}^{MN}.
\]

In implementation, Euler discretization is used on the regular grid:
\[
Y_{k+1} = Y_k + \sum_{i=1}^{d} V_i(Y_k)\,\Delta X_k^i.
\]

\subsection{Learning problems}
\paragraph{Reconstruction target.}
Target:
\[
z_t^{\mathrm{recon}} = \mathrm{Sig}^{\le m}(X^{(t)}),
\]
where \(\mathrm{Sig}^{\le m}\) is truncated signature (order \(m\), in this project \(m=3\)).

\paragraph{Prediction target.}
Target:
\[
z_t^{\mathrm{pred}} = r_t^{(h)}.
\]

\paragraph{Estimator.}
For both tasks, ridge regression:
\[
\hat z_t = \beta_0 + \beta^\top \Phi(X^{(t)}),
\]
\[
(\hat\beta_0,\hat\beta) = \arg\min_{\beta_0,\beta}
\sum_t \left(z_t-\beta_0-\beta^\top\Phi_t\right)^2 + \lambda\|\beta\|_2^2.
\]
\paragraph{Why ridge is preferred.}
Ridge regularization is used because feature dimension can be large relative to effective sample size, especially in rolling-window settings. The \(\ell_2\) penalty stabilizes coefficients and reduces sensitivity to collinearity, which is common in path-derived features.

\subsection{Evaluation metrics}
\paragraph{Statistical metrics.}
Reconstruction quality is measured with component-wise \(R^2\) and mean \(R^2\). Predictive quality is measured with
\[
\mathrm{IC}=\mathrm{corr}(\hat r_t,r_t),\qquad
R^2 = 1 - \frac{\sum_t (r_t-\hat r_t)^2}{\sum_t(r_t-\bar r)^2}.
\]

\paragraph{Backtest metrics.}
Trading rule:
\[
\text{pos}_t = \mathrm{sign}(\hat r_t)\in\{-1,0,1\}.
\]
Net log-return per step:
\[
g_t = \text{pos}_t r_t - c_t,\qquad
c_t = f\cdot |\text{pos}_t-\text{pos}_{t-1}|,
\]
where \(f\) is fee in decimal units (e.g., 5 bps \(=0.0005\)).
Equity:
\[
E_t=\exp\left(\sum_{u\le t}g_u\right).
\]
Reported metrics are cumulative return \(E_T-1\), annualized Sharpe, hit rate \(\Pr(\mathrm{sign}(\hat r_t)=\mathrm{sign}(r_t))\), maximum drawdown, and average turnover.
\paragraph{Interpretation guide.}
In this context, IC quantifies directional alignment, while \(R^2\) quantifies magnitude fit. A model can have useful IC even with low \(R^2\), especially in noisy short-horizon markets. Backtest metrics therefore remain necessary to translate statistical fit into economic relevance.

\subsection{Validation protocols}
\paragraph{Static split.}
Initial train segment (\(\texttt{train\_frac}\)) then fixed test segment.

\paragraph{Walk-forward.}
At each test time \(t\), fit on \([1,\dots,t-1]\), predict at \(t\).  
This is stricter and reduces look-ahead risk.

\paragraph{Cost stress.}
Run identical strategy under fee set \(\{5,10,20\}\) bps.

\paragraph{Non-overlapping test returns.}
Use \(\texttt{step}\ge h\) (e.g., \(\texttt{step}=h=3\)) to reduce overlap-induced metric inflation.
\paragraph{Leakage control.}
Walk-forward estimation ensures each prediction is made with information available at that time only. Non-overlapping horizons reduce dependence among adjacent outcomes. Together, these choices reduce optimistic bias that often appears in dense rolling evaluations.

\subsection{Implementation details}
\paragraph{Script responsibilities.}
\texttt{scan\_overlap\_markets.py} performs keyword-overlap market discovery. \texttt{teichmann\_rigorous\_pipeline.py} performs feature extraction and reconstruction/prediction evaluation. \texttt{teichmann\_backtest.py} performs backtesting with optional walk-forward refits and fee sweeps. \texttt{sentiment\_io.py} handles sentiment loading and timestamp alignment.

\paragraph{Robustness fix in data retrieval.}
If bounded \texttt{prices-history} queries fail for a token, fallback retrieves full history and slices locally to requested time bounds.

\subsection{Reproducible command set}
\paragraph{Market overlap scan.}
\begin{verbatim}
python3 scan_overlap_markets.py \
  --start 2026-02-08T00:00:00Z \
  --end 2026-02-22T00:00:00Z \
  --keywords "ukraine,russia,putin,zelensky,ceasefire,nato" \
  --max_offsets 8 --limit 500 --include_closed \
  --out_csv out/overlap_scan_ukraine_feb8_feb22_2026.csv
\end{verbatim}

\paragraph{Core model run (single market).}
\begin{verbatim}
python3 teichmann_rigorous_pipeline.py \
  --market_id 540816 \
  --sentiment_csv out/gdelt_ukraine_intensity_1h_feb8_feb22.csv \
  --sentiment_ts_col ts --sentiment_value_col sentiment \
  --dt_s 3600 --min_points 200 \
  --window 48 --lag 3 --horizon 3 --step 1 \
  --sig_order 3 --n_list 8,12 --depth_two --activation exp \
  --n_draws 2 --out_prefix teichmann_ukr_feb
\end{verbatim}

\paragraph{Robust walk-forward backtest (top 5 markets).}
\begin{verbatim}
python3 teichmann_backtest.py \
  --market_ids 540816,665224,567689,693519,681144 \
  --sentiment_csv out/gdelt_ukraine_intensity_1h_feb8_feb22.csv \
  --sentiment_ts_col ts --sentiment_value_col sentiment \
  --dt_s 3600 --window 48 --lag 3 --horizon 3 --step 3 \
  --n_list 8,12 --depth_two --activation exp --n_draws 2 \
  --walk_forward --fee_bps_list 5,10,20 \
  --out_prefix teichmann_ukr_feb_bt_top5
\end{verbatim}

\subsection{Empirical artifacts}
\paragraph{Diagnostic figures.}
\begin{figure}[H]
    \centering
    \IfFileExists{out/results_ukr_reconstruction_comparison.png}{
        \includegraphics[width=\textwidth]{out/results_ukr_reconstruction_comparison.png}
    }{
        \fbox{\parbox{0.9\textwidth}{Missing file: \texttt{out/results\_ukr\_reconstruction\_comparison.png}}}
    }
    \caption{Reconstruction quality (mean \(R^2\)) vs hidden dimension \(N\), depth-1 vs depth-2.}
\end{figure}

\begin{figure}[H]
    \centering
    \IfFileExists{out/results_ukr_predictive_ic_comparison.png}{
        \includegraphics[width=\textwidth]{out/results_ukr_predictive_ic_comparison.png}
    }{
        \fbox{\parbox{0.9\textwidth}{Missing file: \texttt{out/results\_ukr\_predictive\_ic\_comparison.png}}}
    }
    \caption{Predictive IC comparison across model variants.}
\end{figure}

\begin{figure}[H]
    \centering
    \IfFileExists{out/results_ukr_predictive_r2_comparison.png}{
        \includegraphics[width=\textwidth]{out/results_ukr_predictive_r2_comparison.png}
    }{
        \fbox{\parbox{0.9\textwidth}{Missing file: \texttt{out/results\_ukr\_predictive\_r2\_comparison.png}}}
    }
    \caption{Predictive \(R^2\) comparison across model variants.}
\end{figure}

\begin{figure}[H]
    \centering
    \IfFileExists{out/teichmann_ukr_feb_bt_top5_mkt540816_best_equity.png}{
        \includegraphics[width=0.9\textwidth]{out/teichmann_ukr_feb_bt_top5_mkt540816_best_equity.png}
    }{
        \fbox{\parbox{0.9\textwidth}{Missing file: \texttt{out/teichmann\_ukr\_feb\_bt\_top5\_mkt540816\_best\_equity.png}}}
    }
    \caption{Example walk-forward equity curve under fee constraints.}
\end{figure}

\paragraph{Aggregate robustness table.}
\IfFileExists{out/teichmann_ukr_feb_bt_top5_aggregate.csv}{
\pgfplotstabletypeset[
    col sep=comma,
    string type,
    columns={market_id,fee_bps,best_cum_return,best_sharpe,median_cum_return,median_sharpe},
    columns/market_id/.style={column name=Market ID},
    columns/fee_bps/.style={column name=Fee (bps)},
    columns/best_cum_return/.style={column name=Best CumRet, fixed, precision=4},
    columns/best_sharpe/.style={column name=Best Sharpe, fixed, precision=3},
    columns/median_cum_return/.style={column name=Median CumRet, fixed, precision=4},
    columns/median_sharpe/.style={column name=Median Sharpe, fixed, precision=3},
    every head row/.style={before row=\toprule,after row=\midrule},
    every last row/.style={after row=\bottomrule}
]{out/teichmann_ukr_feb_bt_top5_aggregate.csv}
}{
\noindent\textbf{Missing file:} \texttt{out/teichmann\_ukr\_feb\_bt\_top5\_aggregate.csv}
}

\subsection{Results}
\paragraph{Core Ukraine modeling results.}
Using the February 8--22, 2026 Ukraine sentiment window at hourly sampling, depth-two exponential randomized signatures improved reconstruction quality relative to depth-one in both benchmark markets. For market 540816, mean reconstruction \(R^2\) increased from \(-0.063\) and \(-0.023\) (depth-one, \(N=8,12\)) to \(0.066\) and \(0.109\) (depth-two, \(N=8,12\)). For market 610256, mean reconstruction \(R^2\) moved from \(-0.016\) and \(0.037\) (depth-one) to \(0.155\) and \(0.157\) (depth-two). Predictive IC in these two-market runs was positive, with values near \(0.22\) for several configurations, while predictive \(R^2\) remained market-dependent.

\paragraph{Walk-forward fee-stress backtest results (top 5 markets).}
In the top-5 overlap basket and non-overlapping walk-forward protocol (\(\texttt{step}=h=3\)), the mean best cumulative return across markets was \(+2.878\%\) at 5 bps, \(+2.597\%\) at 10 bps, and \(+2.039\%\) at 20 bps. The median best cumulative return was \(+3.210\%\), \(+2.541\%\), and \(+1.217\%\), respectively. Four out of five markets remained positive under each fee regime. At 5 bps, the strongest individual best-config outcomes included market 540816 with cumulative return \(+6.617\%\) and Sharpe \(16.40\), and market 693519 with cumulative return \(+6.504\%\) and Sharpe \(22.03\). One market (681144) remained negative across fee levels, indicating cross-market heterogeneity rather than uniform edge.
\paragraph{Practical meaning of these results.}
The evidence is consistent with a weak-to-moderate sentiment-linked signal that is not universal across contracts. The method appears to identify market subsets where sentiment contributes useful timing information, while other contracts remain dominated by noise or unmodeled drivers. This pattern is expected in event markets and motivates contract-level filtering rather than a single pooled strategy.
\paragraph{Result quality cautions.}
Sharpe values from short evaluation windows should be interpreted as ranking statistics, not final performance estimates. The reported fee-stress robustness is encouraging, but it should be viewed as an intermediate validation stage before any production use.

\subsection{Limitations and acceptance criteria}
\paragraph{Limitations.}
Short effective out-of-sample windows can overstate Sharpe, the microstructure model is simplified to proportional turnover fees, and regime transferability is not guaranteed across contracts and periods.

\paragraph{Acceptance criteria for further deployment tests.}
Advance only when IC remains positive across multiple markets, net returns stay positive under 10--20 bps fee stress, and aggregate performance is not dominated by a single contract.
\paragraph{External validity and future protocol.}
For publication-grade confidence, the next step is a temporally separated holdout regime with pre-specified hyperparameters and no post-hoc market selection. A rolling multi-month evaluation with fixed decision rules would provide stronger evidence on persistence and transportability.

\subsection{Reproducibility statement}
\paragraph{Computational reproducibility.}
All outputs are file-backed in \texttt{out/} and \texttt{data/api/}, with explicit command-line configurations documented above. The analysis is designed to be rerunnable end-to-end from raw sentiment time series and market metadata to final aggregate tables and figures without manual intervention.
\paragraph{Data provenance.}
Market metadata are retrieved from Gamma endpoints and price histories from Polymarket CLOB history endpoints. Sentiment series are produced by project scripts from external data sources and then aligned on a common time grid.

\subsection{Compilation}
\begin{verbatim}
pdflatex report_polymarket_sentiment.tex
\end{verbatim}
Run twice if cross-references are introduced later.

\end{document}

